\section{Conclusion and Future Work}

\subsection{Conclusion}
This thesis presented a complete AI/software system for citation-grounded Vietnamese labor-law question answering. The system combines hierarchy-aware corpus processing, Qdrant hybrid retrieval, Neo4j graph expansion, deterministic citation validation, FastAPI backend services, a Next.js frontend, and deterministic evaluation. The work is intentionally scoped: it supports legal information access over six indexed document groups and refuses unsupported questions when the indexed corpus is insufficient.

The main design lesson is that legal QA requires more than semantic similarity. Vietnamese labor-law answers must preserve hierarchy, exact statutory coordinates, cross-references, and legal rank. A flat or hybrid retriever can recover many relevant passages, but it can miss implementing regulations, appendices, or procedure provisions. Graph expansion helps by traversing legal structure around retrieved seeds, while deterministic citation validation prevents final answers from citing outside retrieved evidence.

The final 100-query benchmark shows both strengths and limitations. Graph-augmented retrieval improves Recall@10 from 0.735 to 0.835 and retrieval pass rate from 0.649 to 0.766. The final end-to-end pass rate is 0.780. Citation validation pass rate is 1.000, and out-of-corpus QA pass rate is 1.000. The main failures remain retrieval missing required context, over-expansion to distracting citations, low-information extractive answers, and one answer missing a required rule.

\subsection{Future Work}
Future work should first improve query understanding for informal, paraphrased user questions. Many failures occur when the user describes a practical situation without naming the legal issue. A safer query rewriting layer could map these questions to legal topics and article candidates before retrieval.

Second, table and appendix retrieval should become a first-class subsystem. Retirement schedules, permitted-work lists, and official templates are often stored as table or appendix evidence. Treating these as ordinary text chunks loses important structure. Future work should parse tables into row-level records with explicit coordinate metadata.

Third, graph expansion needs stronger hard-negative filtering. The final evaluation shows that graph expansion improves recall but can introduce related yet distracting legal bases. Future expansion should consider edge type, query intent, legal rank, and benchmark-learned negative patterns before promoting related provisions.

Fourth, labor-dispute procedure retrieval should be improved. These questions require cross-code links between labor rights and civil procedure rules. A targeted procedure module may be more reliable than broad graph expansion.

Fifth, the corpus should be expanded and versioned. A practical legal assistant must track amendments, new decrees, new circulars, validity periods, and superseded provisions. Each corpus update should trigger index rebuilds, graph rebuilds, and benchmark reruns.

Finally, evaluation should include expert review. Deterministic software checks are valuable because they are reproducible and expose system failures clearly. They cannot replace legal expert judgment about interpretation, completeness, and practical advice. The next benchmark should combine deterministic citation tests, expert-labeled answer reviews, and anonymized real user queries.

The implemented system shows that a citation-grounded legal assistant is best treated as an integrated software system. Retrieval, graph structure, validation, refusal behavior, deployment, and evaluation all matter. The language model is only one component. For legal QA, the larger engineering challenge is building the evidence pipeline around it.
